\documentclass[10pt,a4paper]{article}
\usepackage[hmargin=1.5cm,top=1.1cm,bottom=1.1cm]{geometry}
\usepackage[T1]{fontenc}
\usepackage[utf8]{inputenc}
\usepackage{lmodern}
\usepackage{microtype}
\usepackage{booktabs}
\usepackage{graphicx}
\usepackage{xcolor}
\usepackage{tcolorbox}
\usepackage{enumitem}
\usepackage[hidelinks]{hyperref}
\usepackage{soul}
\sethlcolor{yellow!55}
\soulregister\texttt7
\soulregister\emph7
\soulregister\textbf7

\setlist{nosep,leftmargin=*}
\setlength{\parskip}{2.5pt}
\setlength{\parindent}{0pt}
\pagestyle{empty}
\newcommand{\figdir}{../output/figures}
\newcommand{\q}[1]{\textbf{\color{blue!50!black}#1}}

\begin{document}
\small

{\large\textbf{Can a perceptron tag a fashion catalogue? MLPs on Fashion-MNIST}}\\
Ignacio Queipo de Llano \textperiodcentered{} Predictive Machine Learning in Finance, IESE MiF \textperiodcentered{} 27 November 2026\\[-6pt]
\rule{\linewidth}{0.4pt}

\begin{tcolorbox}[colback=white,colframe=black,boxrule=0.4pt,sharp corners,left=5pt,right=5pt,top=3pt,bottom=3pt]
\textbf{Bottom line.} A flat perceptron tags 87\% of the test images correctly, and \textbf{the choice that matters
most is not the architecture but dividing the pixels by 255}: the same (2, 2) network goes from 31.6\% to 72.7\%.
After that, width helps until about 16 neurons and a 2-neuron bottleneck hurts. The errors are the ones a person would
make, among tops with the same silhouette, which is where a model that sees shapes (the CNN, this afternoon) should
gain. Binarising the input loses 3 points and trains 2.5 times faster: the trade-off behind quantised LLMs.
\end{tcolorbox}

\begin{minipage}[t]{0.50\linewidth}\vspace{0pt}
\q{Q1. The data.} 60,000 training and 10,000 test images of $28\times28$ pixels, values 0--255, 6,000 training
images per class.

\q{Q2. Which classes are hard?} \hl{T-shirt, shirt, pullover and coat share a silhouette}; trousers, bags and footwear
are easy.

\includegraphics[width=\linewidth]{\figdir/q2_one_per_class.png}

\q{Q3. Normalisation} (same (2, 2) network, 60 epochs): \hl{raw pixels 31.6\% test; pixels / 255: 72.7\%.}

\q{Q4. Why.} scikit-learn starts the weights within $\pm0.086$. \hl{On raw pixels a neuron's first sum averages 116;
on pixels / 255, 0.45.} Sums that large push ReLUs into huge outputs or into zero, gradients are enormous for some
weights and nil for others, and Adam's step size was designed for inputs near 1. Same lesson as standardising before
K-Means: put every input on the scale the algorithm assumes.

\q{Q5. Architecture.} \hl{(16): 86.0\%; (16, 2): 82.8\%; (32, 32): 87.2\%.} A bottleneck is a layer narrower than the
problem: whatever 16 neurons learned must be squeezed into 2 numbers to choose among 10 classes. \hl{(16, 2) has the
same parameters as (16) and loses 3 points.}
\end{minipage}\hfill
\begin{minipage}[t]{0.47\linewidth}\vspace{0pt}
\includegraphics[width=\linewidth]{\figdir/q7_confusion_matrix.png}

\q{Q7. Confusions} (true $\to$ predicted, (32, 32)): \hl{T-shirt $\to$ shirt 215}, pullover $\to$ shirt 99, coat
$\to$ pullover 93, pullover $\to$ coat 86. Best: trouser 96.7\%, ankle boot 95.8\%; worst: \hl{T-shirt 74.0\%, shirt
77.4\%}. It matches intuition.
\end{minipage}

\vspace{3pt}
\begin{minipage}[t]{0.38\linewidth}\vspace{0pt}
\q{Q6. Parameters of (32, 32).}

\vspace{2pt}
\begin{tabular}{lrr}
\toprule
Layer & Weights + biases & Total\\
\midrule
$784\to32$ & $784\times32+32$ & 25,120\\
$32\to32$ & $32\times32+32$ & 1,056\\
$32\to10$ & $32\times10+10$ & 330\\
\midrule
 & & \hl{\textbf{26,506}}\\
\bottomrule
\end{tabular}

\vspace{2pt}
\hl{95\% of them are in the first layer}: one weight per pixel per neuron.
\end{minipage}\hfill
\begin{minipage}[t]{0.58\linewidth}\vspace{0pt}
\q{Q8. Comparison} (60 epochs, CPU, one model at a time):

\vspace{2pt}
\begin{tabular}{llrrrr}
\toprule
Hidden & Input & Params & Train & Test & CPU\\
\midrule
(2, 2) & 0--255 & 1,606 & 31.9\% & 31.6\% & 12 s\\
(2, 2) & / 255 & 1,606 & 74.0\% & 72.7\% & 16 s\\
(16) & / 255 & 12,730 & 89.4\% & 86.0\% & 28 s\\
(16, 2) & / 255 & 12,624 & 85.9\% & 82.8\% & 47 s\\
(32, 32) & / 255 & 26,506 & 92.5\% & \hl{\textbf{87.2\%}} & 122 s\\
(32, 32) & binary & 26,506 & 92.9\% & 84.0\% & 48 s\\
\bottomrule
\end{tabular}
\end{minipage}

\vspace{3pt}
\hl{More parameters buy accuracy with diminishing returns and a growing train--test gap} (92.5 vs 87.2: mild
overfitting); going from (2, 2) to (16) is worth 13 points, from (16) to (32, 32) barely one, for four times the time.

\q{Q9. Binarise} (pixel $> 0.3 \to 1$). \hl{Test 87.2\% $\to$ 84.0\%; training 122 s $\to$ 48 s.} The model keeps the
shape and loses the grey shades (texture, prints). Quantisation does the same to a model's weights: LLMs stored in 8
or 4 bits instead of 16 or 32 run on cheaper hardware at a small cost in quality.
\end{document}
